% !PRINTED 179
\weberpage{179}
%<<CONT
then $\nu$ is the degree of the group $T_r$, and if the function $\Theta$, which
we applied in \S.\ 43 to the solution of the form problem,
passes, under the substitutions of $T_r$, into
\[ \Theta,\ \Theta_1,\ \ldots \Theta_{\nu-1} \]
then
\[ \Phi_r(t) = (t - \Theta)(t - \Theta_1) \ldots (t - \Theta_{\nu-1}) \]
is a function of $t$ in the extended domain of rationality
$\Omega_r$; the field $\Omega(\Theta)$ is identical with $\Omega_r(\Theta)$. It is therefore an algebraic field of degree $\mu$ over $\Omega$ and of degree $\nu$ over $\Omega_r$.
The form problem of degree $\mu$ is accordingly reduced, by the adjunction
of a radical, to a form problem of degree $\nu$.

\begin{center}
\S.\ 46.

\gesperrt{The extended concept of invariant.}
\end{center}

The relative invariants are at bottom a special case
of a more general concept, which we can apply to the reduction
of the form problem in a similar way to the relative
invariants.

We seek systems of homogeneous forms of equal
degree in the variables $x_1, x_2 \ldots x_n$:
\weq{1}{X_1, X_2, \ldots X_m,}
of such a nature that, under the applications of the substitutions of the group $S$, the functions $X_1, X_2, \ldots X_m$ undergo a
system $\Sigma$ of \gesperrt{linear substitutions}; the substitutions $\Sigma$ then likewise form a group, and indeed
a group which is simply or multiply isomorphic with $S$. If
$m = 1$, we obtain, as one sees, the concept of the relative
invariants. We now seek all the substitutions of $S$ which
leave each single one of the functions (1) unchanged. These substitutions form a group, which we denote by $T$, and
of which we wish to show that it is a normal divisor of
$S$. If we denote by $s$ a substitution of $S$ under which
the $X$ undergo the substitution $\sigma$, then under $s^{-1}$ the $X$ undergo
the substitution $\sigma^{-1}$. If further $\tau$ is a substitution of $T$,
then under $\tau$ the $X$ remain unchanged. Accordingly, under
$s^{-1}\tau s$ the $X$ undergo the substitution $\sigma^{-1}\sigma = 1$, i.\,e. they remain un-
%CONT>>
